\begin{flushright}
  \fechaclase{02 de septiembre de 2026}
\end{flushright}

\paragraph{Teorema de Poincaré--Bendixson}
Considere un sistema dinámico plano y conjunto cerrado y acotado
$m\subset\mathbb{R}^{n}$ tal que

\begin{itemize}
  \item $m$ no contiene puntos de equilibrio o contiene un solo equilibrio
  inestable.
  \item Toda trayectoria que inicia en $m$ permanece en $m$ para todo
  $t\geq 0$, es decir, $m$ permanece alrededor del borde.
\end{itemize}

Entonces $m$ contiene al menos una órbita periódica.

\paragraph{Teorema de estabilidad de Lyapunov}
Sea $X^{*}=0$ un punto de equilibrio para el conjunto
$D\subset\mathbb{R}^{n}$ que contiene $X^{*}=0$. Sea
$V\colon D\to\mathbb{R}$ continuamente diferenciable tal que:
\[
  V(X^{*})=0
\]
\[
  V(X)>0, \qquad \text{en }D-\{0\}
\]
\[
  \dot{V}(X)\leq 0, \qquad \text{en }D
\]
entonces $X^{*}=0$ es estable, además si
\[
  \dot{V}(X)<0, \qquad \text{en }D-\{0\}
\]
entonces $X^{*}$ es asintóticamente estable.

\textit{Nota: $D$ es un espacio delimitado.}

\paragraph{Teorema de Barbashin--Krasovskii}
Sea $X^{*}=0$ un punto de equilibrio y
$V\colon\mathbb{R}^{n}\to\mathbb{R}$ continuamente diferenciable tal que
\[
  V(0)=0
  \qquad\text{y}\qquad
  V(X)>0 \quad \forall X\neq 0
\]
\[
  \lVert X\rVert\to\infty
  \quad\Rightarrow\quad
  V(X)\Rightarrow\infty
\]
\[
  \dot{V}(X)<0 \quad \forall X\neq 0
\]
entonces $X^{*}=0$ es globalmente asintóticamente estable.

\paragraph{Teorema}
Sea $X^{*}=0$ un punto de equilibrio y sea
$V\colon D\to\mathbb{R}$ continuamente diferenciable tal que
\[
  V(0)=0
  \qquad\text{y}\qquad
  V(X)>0
\]
para algún $X_0\in D$ con $\lVert X_0\rVert$ arbitrariamente pequeña, si
\[
  \dot{V}(X_0)>0
\]
entonces el equilibrio es inestable.